\subsection{Commit Analyzer Plug-ins}
\label{sec:CommitAnalyzerPlugins}
\ToDo{
\begin{itemize}
	\item Definition and usage of plug-in-specific configuration parameters in introduction of parent section?
	\item Purpose of plug-in type
	\item Typical input
	\item Realization via implementation of abstract class along example
	\begin{itemize}
		\item Processing data model element instances
		\item Storing results
	\end{itemize}
\end{itemize}
}

\begin{figure}[ht] % t für top, b = bottom, h = here
	\centering
		%trim={<left> <lower> <right> <upper>}
		\includegraphics[width=\columnwidth,trim={1,7cm 20,1cm 5,8cm 3,3cm},clip]{inserts/abstract_commit_analyzer.pdf}
  \caption{\thetool{} \texttt{AbstractCommitAnalyzer} class}
	\label{fig:AbstractCommitAnalyzerClass}
\end{figure}

\begin{figure}[ht]
	\centering
		\lstinputlisting[caption={Blueprint of a \thetool{} commit analyzer main class},label=lst:CommitAnalyzerBlueprint,basicstyle=\small,language=Java]{inserts/CommitAnalyzer.java}
\end{figure}